\documentclass[10pt, letterpaper]{article}
\usepackage[utf8]{inputenc}
\usepackage[margin=0.5in]{geometry}
\usepackage{titlesec}
\usepackage{enumitem}
\usepackage{hyperref}
\usepackage{helvet}
\renewcommand{\familydefault}{\sfdefault}

% Section formatting
\titleformat{\section}{\large\bfseries}{}{0em}{}[\titlerule]
\titlespacing{\section}{0pt}{5pt}{3pt}
\setlist[itemize]{noitemsep, topsep=0pt, leftmargin=1.2em, partopsep=0pt, parsep=0pt}
\pagestyle{empty}

\begin{document}

\begin{center}
    {\huge \textbf{Benjamin Rheault}} \\
    \vspace{2pt}
    (786)-728-7149 | {brheault@ufl.edu} | Gainesville, FL
\end{center}

\section{EDUCATION}
\noindent \textbf{University of Florida - Herbert Wertheim College of Engineering} \hfill Gainesville, FL \\
\noindent \textit{PhD in Computer Science} (GPA: 3.80) \hfill Expected August 2027 \\
\noindent \textit{Master of Science in Computer Science} (GPA: 3.80) \hfill May 2025 \\
\noindent \textit{Bachelor of Science in Computer Science, Minor in Digital Arts} (GPA: 3.96) \hfill May 2022 \\
\noindent \textbullet{} \textit{Honors: Cum Laude}

\vspace{4pt}
\noindent \textbf{Southwestern University of Finance and Economics} \hfill Chengdu, China \\
\textit{Study Abroad Program} \hfill Summer 2019

\section{WORK EXPERIENCE}
\noindent \textbf{Graduate Teaching Assistant}, \textit{University of Florida} \hfill Jan 2025 -- Present \\
\textit{Databases, Discrete Structures, Programming Language Concepts}
\begin{itemize}
    \item Developed a supplementary lecture series covering NoSQL, Cloud DBs, Sharding, and Database Security.
    \item Managed a team of undergraduate TAs; overseeing grading, office hours, and curriculum support for high-enrollment courses.
\end{itemize}

\vspace{0.2cm}

\noindent \textbf{Graduate Research Assistant}, \textit{University of Florida} \hfill Aug 2022 -- Dec 2024
\begin{itemize}
    \item \textbf{PTG UI (DARPA-funded):} Developed an AR user interface for an AI-assisted guidance system. Built ROS communication between HoloLens and PC-based computer vision models to provide real-time procedural task updates.
    \item \textbf{User Modeling:} Led a 50-participant study collecting biometric data under stress. Trained machine learning models (Neural Networks, Random Forests) achieving \textbf{98\% accuracy} in predicting stress states.
    \item \textbf{Blindspots:} Developed a VR screening system for stroke recovery patients, creating a portable digital alternative to the gold-standard Humphrey machine.
    \item \textbf{Rendering Study:} Investigated the impact of CT scan rendering styles on structural comprehension, comparing VR environments to traditional 2D monitors.
\end{itemize}

\vspace{0.2cm}

\noindent \textbf{Undergraduate Teaching Assistant}, \textit{University of Florida} \hfill Aug 2019 -- May 2021 \\
\textit{Computer Organization}
\begin{itemize}
    \item Developed and delivered weekly lectures to simplify complex low-level computer architecture concepts for students. Graded coursework and proctored examinations.
\end{itemize}

\section{PUBLICATIONS}
\begin{itemize}
    \item\textit{Data Dump to Digestible Chunks: Automated Summarization of Provenance Logs for Communication} [ArXiv]
    \item \textit{Effects of Technique Awareness on the Detection of Remapped Hands in Virtual Reality.} IEEE VR 2024. (Paper)
    \item \textit{Predictive Task Guidance with Artificial Intelligence in Augmented Reality.} IEEE VR 2024. (Poster)
    \item \textit{Multitasking with Graphical Encoding Visualization of Numerical Values in Virtual Reality.} IEEE VR 2024. (Poster)
    \item \textit{Pseudocode vs. Compiler: Comparing Measures of Student Programming Ability in CS1 and CS2.} ITiCSE 2023. (Paper)
    \item \textit{Analysis of Effect of Answering Reflection Prompts in a Computer Organization Class} ASEE 2022. (Paper)
\end{itemize}

\section{PROJECTS}
\noindent \textbf{Field Particle (Python, PyQt5)} \hfill Jan 2025 -- May 2025
\begin{itemize}
    \item Collaborated with Mechanical Engineering faculty to build a particle collision detection simulator for computational fluid dynamics (CFD) research.
\end{itemize}

\noindent \textbf{ONE.UF Mobile (TypeScript)} \hfill Jan 2022 -- May 2022
\begin{itemize}
    \item Developed a cross-platform mobile application version of the University of Florida's official student portal.
\end{itemize}

% \noindent \textbf{Production Graph (Python, Django, Rust)} \hfill May 2020 -- Aug 2020
% \begin{itemize}
%     \item Built a web service to estimate labor hours and commodity costs, validating results against real-world data.
% \end{itemize}

\noindent \textbf{RTS Bus Tracker (JavaScript, Node.js)} \hfill Feb 2020
\begin{itemize}
    \item Developed a Google Assistant action integrated with regional transport APIs; awarded the \textbf{Google Cloud Award} at Swamphacks 2020.
\end{itemize}

\noindent \textbf{Personal Trainer Website (JavaScript, MERN stack)} \hfill Jan 2020 -- May 2020
\begin{itemize}
    \item Developed a web application for a local personal trainer business, allowing clients to create accounts, schedule sessions, and keep track of their fitness goals.
\end{itemize}

\section{SKILLS \& COURSEWORK}
\noindent \textbf{Languages:} Python, Java, C++, C\#, JavaScript, SQL, Pony \\
\noindent \textbf{Platforms/Tools:} Unity, Maya, Blender, ROS (Robot Operating System), Git \\
\noindent \textbf{Relevant Coursework:} Computer Graphics, Machine Learning Engineering, Software Engineering, Data Science (all graduate level)

\end{document}